% ===================================================================
%  TÁBLA Uimh. 5 — An Teach agus conas a thógtar é.
%  Traduction 2026 d'apres texto/io, controlee sur texto/fr, texto/en
%  et texto/es. Consigne : voir texto/ga/10-tabla-01.tex.
%
%  Le tableau 5 est un DIALOGUE, et les deux volumes composent le nom
%  du parleur en petites capitales : on garde \textsc{}.
%
%  Les renvois du plan sont des LETTRES — (a) le couloir, (b) le
%  salon... — et non des numeros, y compris la ou l'imprimeur a
%  compose « (1) » pour « (l) ».
%
%  LE TABLEAU DES METIERS. L'irlandais compose presque tous les noms
%  d'artisan avec le suffixe « -éir » ou « -adóir » : « saor cloiche »
%  le tailleur de pierre, « siúinéir » le menuisier, « pluiméir » le
%  plombier, « gabha » le forgeron. Le dernier est le mot celtique
%  ancien, les autres sont des composes ou des emprunts — et cela suit
%  le critere degage en roumain : le forgeron etait au village depuis
%  toujours, le plombier est arrive avec la plomberie.
% ===================================================================

%%K t05-apar-1 apar
\VUsaut{6.0mm}
\VUtitre{15.0pt}{60}{TÁBLA Uimh. 5}
\VUsaut{1.4mm}
\VUtitre{11.4pt}{0}{An Teach agus conas a thógtar é.}
\VUsaut{2.2mm}

%%K t05-tit-1 sub
\VUpk{10.2pt}{40}{Casadh sona.}
\VUsaut{3.0mm}

%%K t05-01-1 p
\textsc{Eoin}. --- A uncail, ba mhaith liom go mór a fháil amach conas a thógtar
\VUgras{teach}.

%%K t05-01-2 p
\textsc{An tUncail} \textsuperscript{(80)}. --- Níl sé deacair, a bhuachaill; tar
liom agus taispeánfaidh mé duit an ceann atá á thógáil ag an Uasal Bernard
\textsuperscript{(79)} agus ina mbeidh mé i mo chónaí.

%%K t05-01-3 p
\textsc{Eoin}. --- Agus cé a tharraing \VUgras{plean} \textsuperscript{(76)} an tí
seo?

%%K t05-01-4 p
\textsc{An tUncail}. --- An \VUgras{t-ailtire} \textsuperscript{(75)}; rinne sé
líníocht an-soiléir, ceadaíodh í, agus thosaigh an \VUgras{conraitheoir}
\textsuperscript{(77)} ar na hoibreacha.

%%K t05-01-5 p
\textsc{Eoin}. --- Agus cé hé an conraitheoir?

%%K t05-01-6 p
\textsc{An tUncail}. --- An tUasal Blanc, athair do chomrádaí Séamas. Stiúrann sé
na hoibrithe in éineacht leis an Uasal Roux, an t-ailtire.

%%K t05-01-7 p
Seo chugainn é! Tá an tUasal Blanc ag teacht díreach in am lena \VUgras{riail}
\textsuperscript{(78)}, agus an tUasal Roux lena phlean.

%%K t05-01-8 p
Lá maith, a dhaoine uaisle. Conas atá sibh?

%%K t05-01-9 p
\textsc{An tUasal Roux}. --- Go breá, go raibh maith agat. Lá maith, a Eoin; nach
mór atá an buachaill tar éis fás!

%%K t05-01-10 p
\textsc{An tUncail}. --- Go deimhin, fear beag ar fad. Tá sé an-dáiríre; ní mór
gach a bhfeiceann sé a mhíniú dó. Inniu ba mhaith leis a fháil amach conas a
thógtar teach.

%%K t05-tit-2 sub
\VUpk{10.2pt}{40}{Na hOibrithe.}
\VUsaut{3.0mm}

%%K t05-01-11 p
\textsc{An tUasal Blanc}. --- Tar liom mar sin, a chara óig, agus míneoidh mé gach
rud duit láithreach, mar is breá liom páistí ar mian leo foghlaim.

%%K t05-01-12 p
An bhfeiceann tú an t-oibrí sin leis an \VUgras{sluasaid}
\textsuperscript{(82)} ina lámha: is é an \VUgras{tochailteoir}
\textsuperscript{(81)} é. Tá sé ag tochailt \VUgras{díge} \textsuperscript{(46)},
chun an phíopa uisce a chur síos. Is é freisin a thochail an talamh mar a bhfuil an
teach, chun go bhféadfadh an saor cloiche na ceithre \VUgras{bhalla} a thógáil. An
chuid de na ballaí sin atá sa talamh, tugtar \VUgras{bunsraith} uirthi
\textsuperscript{(2)}. Déanann an spás atá dúnta ag an mbunsraith
\VUgras{siléar} \textsuperscript{(3)}. Tá fuinneoigín ag an siléar seo, ar a
dtugtar \VUgras{poll aeir} \textsuperscript{(5)}, \VUgras{doras}
\textsuperscript{(4)} agus staighre.

%%K t05-01-13 p
Lean an \VUgras{saor cloiche} \textsuperscript{(108)} ag tógáil na mballaí, ag
fágáil bearnaí do na \VUgras{doirse} \textsuperscript{(8)} agus do na
\VUgras{fuinneoga} \textsuperscript{(17)}.

%%K t05-01-14 p
\textsc{Eoin}. --- Ach ní fheicim é, an saor cloiche seo!

%%K t05-01-15 p
\textsc{An tUasal Blanc}. --- Tá deireadh curtha aige leis an obair ar an teach
cheana. Tá sé thall ansin, in aice leis an \VUgras{stábla}
\textsuperscript{(54)}, ar a \VUgras{scafall} \textsuperscript{(110)}, ag tógáil
bhalla na \VUgras{níolainne} \textsuperscript{(56)}.

%%K t05-01-16 p
Tá liaigh, umar agus \VUgras{luaidhlíne} aige \textsuperscript{(109)}. Ach ní
chuimhneoidh tú ar na hainmneacha sin.

%%K t05-01-17 p
\textsc{Eoin}. --- Cuimhneoidh go deimhin, mar iarrfaidh mé ar m'uncail liaigh
bheag agus umar, agus déanfaidh mé an luaidhlíne mé féin as sreangán.

%%K t05-tit-3 sub
\VUsaut{3.0mm}
\VUpk{10.2pt}{40}{Na hÁbhair.}
\VUsaut{3.0mm}

%%K t05-01-18 p
\textsc{An tUncail}. --- Go maith. Ach inis dom, a Eoin, cad atá riachtanach chun
teach a thógáil?

%%K t05-01-19 p
\textsc{Eoin}. --- Tá \VUgras{cloch shnoite} \textsuperscript{(59)} agus
\VUgras{moirtéal} \textsuperscript{(65)} riachtanach.

%%K t05-01-20 p
\textsc{An tUncail}. --- Díreach mar sin. Féach ar an bhfear sin leis an hata mór,
ar a \VUgras{chairt} \textsuperscript{(136)}. Is é an \VUgras{cairteoir}
\textsuperscript{(135)} é. Gach lá tugann sé anseo \VUgras{bloic chloiche}
\textsuperscript{(60)}. Díluchtaíonn sé an chairt sa chlós, agus snoíonn na
\VUgras{saoir chloiche} \textsuperscript{(83)} na bloic agus gearrann siad lena
\VUgras{sábh cloiche} \textsuperscript{(85)} iad. Iompraíonn an \VUgras{fear
cabhrach} \textsuperscript{(146)} chuig an saor cloiche iad, agus cuireann seisean
ina n-áit iad.

%%K t05-01-21 p
Idir an dá linn ullmhaíonn an \VUgras{meascthóir} \textsuperscript{(87)} an
moirtéal. Tá \VUgras{gaineamh} \textsuperscript{(62)}, \VUgras{aol}
\textsuperscript{(63)} agus \VUgras{uisce} \textsuperscript{(64)} de dhíth air. Tá
an t-aol ina aice i \VUgras{mála} \textsuperscript{(71)}; tugann \VUgras{cúntóir an
tsaoir} \textsuperscript{(111)} gaineamh chuige i \VUgras{bhara}
\textsuperscript{(112)}, agus tá an \VUgras{printíseach} \textsuperscript{(113)} ag
an gcaidéal \textsuperscript{(58)}. Líonann sé an \VUgras{buicéad}
\textsuperscript{(114)}. Beidh an moirtéal réidh gan mhoill. Tógann an
\VUgras{tabharthóir} \textsuperscript{(107)} é agus iompraíonn sé ar an staighre é
go dtí an scafall, mar a bhfuil an saor cloiche ag críochnú an taoibh seo den
teach.

%%K t05-01-22 p
Agus anois: cad a thugtar ar an oibrí a dhéanann an \VUgras{díon}
\textsuperscript{(31)}?

%%K t05-01-23 p
\textsc{Eoin}. --- Is é an \VUgras{slinnteoir} \textsuperscript{(118)} é.

%%K t05-tit-4 sub
\VUsaut{3.0mm}
\VUpk{10.2pt}{40}{Díon an tí.}
\VUsaut{3.0mm}

%%K t05-01-24 p
\textsc{An tUasal Blanc}. --- Sea, ach ar dtús tagann an \VUgras{siúinéir tí}
\textsuperscript{(115)}.

%%K t05-01-25 p
Déanann an siúinéir tí na \VUgras{creataí} \textsuperscript{(35)}. Nascann sé na
\VUgras{bíomaí} \textsuperscript{(37)} le \VUgras{dúigh} \textsuperscript{(117)}.
Is siúinéirí tí iad na hoibrithe sin thuas, ina mbrístí leathana veilbhite.
Coinníonn duine acu píosa adhmaid, snoíonn an duine eile dúigh lena
\VUgras{thua} \textsuperscript{(116)}. Is é an duine atá in aice leis an
\VUgras{simléar} \textsuperscript{(41)} an slinnteoir. Buaileann sé na slata ar na
\VUgras{bíomaí} (*), agus an \VUgras{slinn} \textsuperscript{(66)} ar na slata.
Uaireanta cuireann sé tíleanna.

%%K t05-noto-1 noto
\VUnotes{143.2mm}{%
(*) \VUgras{Balk-o} = Gearm. \textit{Balken}, Béarla \textit{joist}, Fr.
\textit{solive}, Iod. \textit{travicello}, Gaeilge \textit{bíoma}, Sp. agus Port.
\textit{viga}. Fréamh theicniúil, nár glacadh léi fós, ach a mholtar anseo chun
brí an fhocail Fhrancaigh \textit{solive} a thabhairt, focal nach é \textit{trabo}
(Fr. \textit{poutre}) é ná píosa adhmaid. Is lomán snoite é a iompraíonn an
t-urlár agus an tsíleáil.%
}

%%K t05-01-26 p
Is de \VUgras{thíleanna} \textsuperscript{(55)} díon an stábla, is de
\VUgras{shlinn} \textsuperscript{(32)} díon an tí, agus is de \VUgras{shinc}
\textsuperscript{(24)} díon na vearanda.

%%K t05-01-27 p
\textsc{Eoin}. --- An é an slinnteoir a dhéanann na díonta since freisin?

%%K t05-01-28 p
\textsc{An tUasal Blanc}. --- Ní hé, is é an \VUgras{stánadóir}
\textsuperscript{(121)} nó an \VUgras{pluiméir} é, an fear sin atá ag cur an
\VUgras{tsolasfhuinneog} \textsuperscript{(38)} i ndíon an tí. Is é freisin a
chuireann na \VUgras{gáitéir} \textsuperscript{(43)} agus na \VUgras{píopaí uisce}
\textsuperscript{(42)}. Sádrálann sé sinc agus stán. Is é a dhaingnigh an
\VUgras{tintreach-shlat} \textsuperscript{(40)} agus an \VUgras{coileach gaoithe}
\textsuperscript{(39)} ar an \VUgras{bhinn} \textsuperscript{(36)}.

%%K t05-01-29 p
\textsc{Eoin}. --- Mar sin tá an teach réidh?

%%K t05-tit-5 sub
\VUsaut{3.0mm}
\VUpk{10.2pt}{40}{Na hUirlisí.}
\VUsaut{3.0mm}

%%K t05-01-30 p
\textsc{An tUasal Blanc}. --- Ní hé go hiomlán. Tógann an \VUgras{plástrálaí}
\textsuperscript{(99)} as \VUgras{bríce} \textsuperscript{(61)} na landairí, a
roinneann é ina sheomraí éagsúla. Cuireann sé an \VUgras{plástar}
\textsuperscript{(70)} ar na \VUgras{síleálacha} \textsuperscript{(26)} agus ar na
ballaí. Anois tá sé ag obair ar shíleáil na \VUgras{vearanda}
\textsuperscript{(33)}. Mar atá ag an saor cloiche, tá \VUgras{liaigh}
\textsuperscript{(100)} agus \VUgras{umar} \textsuperscript{(101)} aige.

%%K t05-01-31 p
Ansin déanann an siúinéir na doirse, na fuinneoga, na \VUgras{hurláir adhmaid}
\textsuperscript{(16)} agus na \VUgras{comhlaí} \textsuperscript{(19)}.

%%K t05-01-32 p
Anois tá sé ag obair sa chlós ag a \VUgras{bhinse oibre} \textsuperscript{(90)}.
Tá \VUgras{sábh} aige \textsuperscript{(94)}, chun an \VUgras{t-adhmad}
\textsuperscript{(69)} a ghearradh, \VUgras{locar} \textsuperscript{(91)},
\VUgras{teanntán} \textsuperscript{(92)}, \VUgras{siséal} \textsuperscript{(94
\textit{bis})}, \VUgras{compás} \textsuperscript{(95)} agus \VUgras{casúr}
adhmaid \textsuperscript{(95 \textit{bis})}.

%%K t05-01-33 p
\textsc{Eoin}. --- Á, is fearr fós a bheith i do shiúinéir ná i do shaor cloiche! A
uncail, an gceannóidh tú dom binse oibre, sábh agus locar? Déanfaidh mé bothán do
Bhran.

%%K t05-01-34 p
\textsc{An tUncail}. --- Ceannóidh, a bhuachaill.

%%K t05-01-35 p
\textsc{Eoin}. --- Nach áthasach atáim! Agus cé hé an fear sin atá ina sheasamh ag
an doras tosaigh?

%%K t05-tit-6 sub
\VUsaut{3.0mm}
\VUpk{10.2pt}{40}{Obair an phéintéara.}
\VUsaut{3.0mm}

%%K t05-01-36 p
\textsc{An tUasal Blanc}. --- Is é an \VUgras{péintéir} \textsuperscript{(102)} é.
Tá sé ina sheasamh ar a dhréimire le pota \VUgras{péinte}
\textsuperscript{(103)} agus lena \VUgras{scuab} \textsuperscript{(104)}.
Clúdaíonn sé adhmad an dorais tosaigh, chun go mairfidh sé níos faide.

%%K t05-01-37 p
Is é freisin a ghreamaíonn na páipéir bhalla sna seomraí.

%%K t05-01-38 p
Tá an \VUgras{obairtheoir} \textsuperscript{(107)} ar an \VUgras{dréimire beag}
\textsuperscript{(106)}, ar an \VUgras{mbalcóin} \textsuperscript{(28)}, ag crochadh
an \VUgras{dallóige} \textsuperscript{(29)}.

%%K t05-01-39 p
Is é an t-oibrí atá ina sheasamh ag an bhfuinneog mhór an \VUgras{gloineadóir}
\textsuperscript{(96)}. Cuireann sé na \VUgras{gloiní} \textsuperscript{(18)} ar
\VUgras{phutaí} \textsuperscript{(73)}. Agus is é an t-oibrí sin atá ar a ghlúine
ag doras an tsiléir an \VUgras{glasadóir} \textsuperscript{(97)}. Tá sé ag cur an
\VUgras{ghlais} \textsuperscript{(6)}. Tá an \VUgras{liomhán}
\textsuperscript{(98)} ina aice.

%%K t05-01-40 p
\textsc{Eoin}. --- Agus an té atá ag bualadh le \VUgras{casúr}
\textsuperscript{(84)} ar phíosa iarainn, an glasadóir é sin freisin?

%%K t05-01-41 p
\textsc{An tUasal Blanc}. --- Níos cruinne, is \VUgras{gabha}
\textsuperscript{(125)} é. Gaibhníonn sé na \VUgras{páirteanna iarainn}
\textsuperscript{(67)} don \VUgras{leictreoir} \textsuperscript{(124)}, atá ina
sheasamh ar dhíon na \VUgras{cóiste-thí} \textsuperscript{(51)}. Tá
\VUgras{teallach} \textsuperscript{(126)}, \VUgras{inneoin}
\textsuperscript{(127)} agus \VUgras{teanchair} \textsuperscript{(128)} aige.

%%K t05-01-42 p
Daingníonn an leictreoir \VUgras{sreang chopair} \textsuperscript{(44)} an
teileafóin.

%%K t05-tit-7 sub
\VUpk{10.2pt}{40}{Obair sa ghairdín.}
\VUsaut{3.0mm}

%%K t05-01-43 p
\textsc{An tUncail}. --- Ag cúl an \VUgras{chlóis} \textsuperscript{(45)}, in
aice leis na \VUgras{ráillí} \textsuperscript{(47)} agus leis an
\VUgras{ngeata} \textsuperscript{(48)}, feiceann tú an \VUgras{garraíodóir}
\textsuperscript{(129)}. Tá sé ag ullmhú \VUgras{gairdín} bhig
\textsuperscript{(49)}. Rómhair sé an talamh lena \VUgras{spád}
\textsuperscript{(131)}, cuireann sé síolta agus réitíonn sé leis an
\VUgras{ráca} \textsuperscript{(130)}. Ansin as a \VUgras{channa uisce}
\textsuperscript{(132)} uiscíonn sé na \VUgras{ceapacha} \textsuperscript{(150)},
agus fásfaidh na plandaí.

%%K t05-01-44 p
Tá an \VUgras{giolla capaill} \textsuperscript{(133)}, atá díreach tar éis an capall
a ghlanadh agus a bheathú, ag glanadh na \VUgras{cóiste} \textsuperscript{(52)} le
spúinse, le \VUgras{caidéal láimhe} \textsuperscript{(134)} agus le buicéad uisce.
Ansin cuirfidh sé isteach sa chóiste-theach í agus glasálfaidh sé an doras le
\VUgras{bolta} trom \textsuperscript{(53)}.

%%K t05-01-45 p
Sin é, anois tá tú sásta, is dóigh liom; bhí míniú go leor ann.

%%K t05-01-46 p
\textsc{Eoin}. --- Tá, a uncail, ach ba mhaith liom an teach a fheiceáil ón taobh
istigh.

%%K t05-01-47 p
\textsc{An tUncail}. --- Beidh sin amárach. Míneoidh an tUasal Roux an plean duit
agus ainmneoidh sé gach seomra.

%%K t05-tit-8 sub
\VUsaut{3.0mm}
\VUpk{10.2pt}{40}{Leagan amach an tí.}
\VUsaut{2.4mm}

%%K t05-apar-2 apar
{\centering\textit{(Féach an plean.)}\par}
\VUsaut{2.4mm}

%%K t05-01-48 p
\textsc{An tAiltire}. --- Téimis, a Eoin, treoróidh mé thú trí chodanna éagsúla an
tí.

%%K t05-01-49 p
Téimis isteach san \VUgras{urlár íochtarach} \textsuperscript{(7)}. Seo cnaipe an
\VUgras{chloigín} \textsuperscript{(11)}. Ardaímid trí \VUgras{chéim}
\textsuperscript{(14)} an \VUgras{phóirse} \textsuperscript{(9)}, trasnaímid
\VUgras{tairseach} \textsuperscript{(10)} an \VUgras{dorais tosaigh}
\textsuperscript{(8)} --- agus seo sinn ag bun an \VUgras{staighre}
\textsuperscript{(13)}.

%%K t05-01-50 p
\textsc{Eoin}. --- Is é seo an conair.

%%K t05-01-51 p
\textsc{An tAiltire}. --- Ní hé, is é seo an \VUgras{halla}
\textsuperscript{(12)}. Tá an \VUgras{chonair} \textsuperscript{(a)} ar chúl. Ar
dheis --- an \VUgras{seomra suí} \textsuperscript{(b)} agus an \VUgras{seomra
bia} \textsuperscript{(c)}; os comhair an tseomra bia --- an \VUgras{chistin}
\textsuperscript{(d)} agus an \VUgras{phantrach} \textsuperscript{(e)}; ansin, ar
thaobh an aghaidh, an \VUgras{oifig} \textsuperscript{(f)}. Tá an
\VUgras{vearanda} \textsuperscript{(23)} lena \VUgras{doras gloine}
\textsuperscript{(25)} greamaithe den seomra suí.

%%K t05-01-52 p
Anois téimis suas an staighre. Coinnigh greim ar an \VUgras{ráille}
\textsuperscript{(15)} agus comhair na céimeanna. Tá ceithre cinn déag ann. Seo an
\VUgras{léibheann} \textsuperscript{(m)}: táimid ar an gcéad \VUgras{urlár}
\textsuperscript{(27)}.

%%K t05-tit-9 sub
\VUsaut{3.0mm}
\VUpk{10.2pt}{40}{An chéad urlár.}
\VUsaut{3.0mm}

%%K t05-01-53 p
Os comhair an staighre --- an \VUgras{leithreas} \textsuperscript{(20)} agus an
\VUgras{seomra maisiúcháin} \textsuperscript{(j)}. Ar dheis \VUgras{seomra
codlata} álainn \textsuperscript{(g)} le dhá fhuinneog ar \VUgras{aghaidh}
\textsuperscript{(1)} an tí. Ar chlé --- an \VUgras{seomra folctha}
\textsuperscript{(1)} agus an \VUgras{seomra aíonna} \textsuperscript{(i)} le
\VUgras{halcóib} mór \textsuperscript{(h)}. In aice leis an seomra codlata --- an
\VUgras{seomra páistí} \textsuperscript{(k)}. Breathnaíonn sé ar \VUgras{ardán}
leathan \textsuperscript{(21)}. Faoin ardán sin leithreas eile
\textsuperscript{(20)}, agus in aice leis an mballa --- an \VUgras{clog}
\textsuperscript{(22)}, a bhuailtear don dinnéar.

%%K t05-01-54 p
\textsc{Eoin}. --- Téimis suas go dtí an \VUgras{dara hurlár}.

%%K t05-01-55 p
\textsc{An tAiltire}. --- Níl dara hurlár ann. Faoin díon níl ach na
\VUgras{seomraí áiléir} \textsuperscript{(33)} agus an t-áiléar
\textsuperscript{(34)}. Tá an turas thart, is féidir linn dul síos.

%%K t05-01-56 p
\textsc{Eoin}. --- Go raibh maith agat, a dhuine uasail, bhí sé an-suimiúil.
Inseoidh mé do m'athair gach a bhfaca mé.
\VUsaut{4.0mm}
\VUfilet{22mm}
